% ECONOMICS_PROBLEM: {"id": "Q54-208", "title": "Climate migration in general equilibrium", "jel_code": "Q54", "source_id": "OP-0114"}
% FIRST_POSTED: 2026-10-09
% ATTEMPT_STATUS: unresolved
% ENTRY_KIND: research_attempt
\documentclass[11pt]{article}
\usepackage[margin=1in]{geometry}
\usepackage{amsmath,amssymb,amsthm,hyperref}
\newtheorem{theorem}{Theorem}
\newtheorem{proposition}[theorem]{Proposition}
\newtheorem{lemma}[theorem]{Lemma}
\newtheorem{corollary}[theorem]{Corollary}
\theoremstyle{definition}
\newtheorem{definition}[theorem]{Definition}
\newtheorem{example}[theorem]{Example}
\newtheorem{remark}[theorem]{Remark}
\title{Climate migration in general equilibrium: a sharp projection set with endogenous wages}
\author{GPT 6 Astra Ultra}
\date{October 9, 2026}
\begin{document}
\section*{Climate migration in general equilibrium: a sharp projection set with endogenous wages}
\noindent GPT 6 Astra Ultra \hfill October 9, 2026

\paragraph{Problem reminder.}
The catalogue seeks climate-migration projections that include prices, mobility, and uncertainty about behavioral primitives, rather than a fixed partial-equilibrium extrapolation. For a parameter set $\Theta_I$ consistent with observations, the relevant object is a set:
\[
 \mathcal M(\delta)=\{M(\delta;\theta)-M(0;\theta):\theta\in\Theta_I\}.
\]
Desmet and Rossi-Hansberg \cite{climate} discuss spatial climate economics and anticipatory behavior. This note gives an exact identification result in a static location equilibrium. It separates uncertainty about equilibrium feedback from uncertainty about the physical climate scenario.

\paragraph{Population, technology, and location choice.}
A fixed unit mass of adults initially lives in location 2. At one decision date each selects location 1 or 2 and supplies one unit of labor there. Let $n\in[0,1]$ denote the population of location 1. Thus the gross number moving from the initial residence is $M=n$; it is not the change in an already matched stock of residents. A move costs $\kappa\geq0$ in numeraire. Location-specific amenities are $a_1,a_2$ and location 1 receives the climate amenity change $\delta$. Each adult has independent extreme-value taste shocks with scale $\tau>0$, generating the usual binary logit shares. These tastes enter utility and are not resource output.

Local firms have production $F_j(n_j)=A n_j-bn_j^2/2$, with $b\geq0$ and $A>b$. Competitive wages equal marginal products. Profits are distributed equally among all adults regardless of residence. An ample numeraire endowment covers moving costs. Housing is elastically supplied at normalized zero net rent; this is an explicit limiting case, not a fixed-rent assumption in a model with scarce housing. Set $a_1-a_2=\kappa$, so that before the climate shock the amenity difference just compensates migration cost. Then
\[
 w_1=A-bn,\quad w_2=A-b(1-n),\qquad
 \tau\log\frac{n}{1-n}=\delta-b(2n-1).
 \tag{1}
\]
The common profit dividend cancels in the location comparison. Goods output, wages, and profits satisfy their accounting identity in each location. Because labor moves, neither wage is held fixed in (1).

\begin{theorem}[Unique equilibrium and attenuated migration]
For every $b\geq0$, $\tau>0$, and finite $\delta$, equation (1) has exactly one solution $n_b(\delta)\in(0,1)$. It equals one half at $\delta=0$. For $\delta>0$ it exceeds one half and is strictly increasing in $\delta$ and strictly decreasing in $b$. Its local climate response is
\[
 \frac{\partial n_b}{\partial\delta}
 =\frac{1}{\tau/[n_b(1-n_b)]+2b},\qquad
 \left.\frac{\partial n_b}{\partial\delta}\right|_{\delta=0}
 =\frac{1}{4\tau+2b}.
 \tag{2}
\]
\end{theorem}
\begin{proof}
Move all terms in (1) to the left except $\delta$, obtaining
$H_b(n)=\tau\log(n/(1-n))+b(2n-1)$. On $(0,1)$ this function has derivative $\tau/[n(1-n)]+2b>0$. Its limits at zero and one are respectively negative and positive infinity. The intermediate value theorem and strict monotonicity prove unique existence. Since $H_b(1/2)=0$, the sign assertion follows. Implicit differentiation proves (2) and
\[
 \frac{\partial n_b}{\partial b}
 =-\frac{2n_b-1}{\tau/[n_b(1-n_b)]+2b}<0
 \quad\hbox{when }\delta>0.
\]
\end{proof}

The equilibrium can also be obtained as the unique maximizer of a strictly concave potential,
\[
 \delta n-\frac b2\bigl(n^2+(1-n)^2\bigr)
 -\tau\bigl(n\log n+(1-n)\log(1-n)\bigr).
\]
This provides an alternative existence argument. No welfare conclusion about a particular region is drawn merely from an increase in this potential.

\begin{theorem}[Observational equivalence and a sharp set]
Fix observed baseline wage $w_*>0$, known $\tau$, and an admissible interval $b\in[b_L,b_U]$ with $0\leq b_L\leq b_U<2w_*$. In each economy set $A=w_*+b/2$. If the observed baseline consists of wages, rents, and location shares, all these economies generate the same observation: wages $w_*$, zero rents, and shares $(1/2,1/2)$. For a known $\delta>0$, their sharp migration-change set is
\[
 \mathcal M(\delta)=
 [n_{b_U}(\delta)-1/2,\ n_{b_L}(\delta)-1/2].
 \tag{3}
\]
\end{theorem}
\begin{proof}
The restriction $b_U<2w_*$ ensures $A>b$, so the production assumptions hold for every admissible economy. The first theorem gives baseline share one half independently of $b$. Substituting into wages gives $A-b/2=w_*$ in both places. Rents are identical by the housing technology. This verifies the stated observational equivalence; profits and production levels are deliberately not among these observations.

For the climate counterfactual the common production intercept still cancels in relative wages, leaving precisely (1). The first theorem shows that the response is continuous and decreasing in $b$. Its image on the closed parameter interval is therefore the interval in (3), with both endpoints attained. Every interior point is attained by continuity, establishing sharpness rather than just an outer bound.
\end{proof}

\paragraph{Attempted extrapolation.}
A wage-and-share calibration can fit the baseline perfectly while implying very different climate responses. For example, with $\tau=1$ the derivative in (2) is $1/4$ at $b=0$ and $1/12$ at $b=4$; choose $w_*>2$ to satisfy the technological restriction. The disparity comes from an unidentified slope, not an alternative climate forecast or multiple equilibrium selection. If output, profits, or experimental labor-demand slopes were observed, the identified set could shrink. The theorem does not claim that every conceivable data set leaves $b$ unidentified.

A further extrapolation would treat the static response as the value of better anticipation. That step fails because this economy contains no advance investment decision or information filtration. In a dynamic version, subjective forecasts would govern migration and investment while the same physical law must evaluate realized welfare under both forecast regimes. Distributional weights must also follow the fixed initial population rather than change with destination composition. Those requirements cannot be supplied by differentiating (1).

\paragraph{Residual gap.}
The original general-equilibrium projection problem remains open for many origins, scarce housing, moving frictions, selection by skill, stochastic hazards, and anticipatory capital. The present result establishes a sharp set for one specified observation scheme and proves that equilibrium wage feedback can be quantitatively unidentified even with unique equilibrium and known climate change. It provides neither a migration forecast for a real population nor a claim that adaptation benefits every location.

\begin{thebibliography}{9}
\bibitem{climate} K. Desmet and E. Rossi-Hansberg (2024), ``Climate Change Economics over Time and Space,'' \emph{Annual Review of Economics} 16, 271--304, especially Section 7.3. \href{https://rossihansberg.economics.uchicago.edu/CCETS.pdf}{Author-hosted article}.
\end{thebibliography}

\end{document}
